% Generated draft: requires mathematical and source review.
\documentclass[11pt]{article}
\usepackage[margin=25mm]{geometry}
\usepackage{fontspec}
\setmainfont{Latin Modern Roman}
\usepackage{amsmath,amssymb,mathtools,unicode-math}
\setmathfont{Latin Modern Math}
\usepackage{enumitem,needspace,xurl,hyperref}
\setlength{\parindent}{0pt}
\setlength{\parskip}{5pt}
\setlength{\emergencystretch}{3em}
\newcommand{\N}{\mathbb{N}}
\newcommand{\Z}{\mathbb{Z}}
\newcommand{\Q}{\mathbb{Q}}
\newcommand{\R}{\mathbb{R}}
\newcommand{\C}{\mathbb{C}}
\newcommand{\F}{\mathbb{F}}
\newcommand{\A}{\mathbb{A}}
\newcommand{\PP}{\mathbb{P}}
\newcommand{\E}{\mathbb{E}}
\newcommand{\eps}{\varepsilon}
\newcommand{\dd}{\,\mathrm{d}}
\begin{document}
\section*{UnsolvedMath ID: 30001417}
\textbf{Source problem code:} catalogue entry
\subsection{Definitions and mathematical statement}
\textbf{Endpoint Korn--Sobolev question.} The supplied record asks for the optimal constant in
\[
\left(\int_{\R^n}|u(x)|^{n/(n-1)}\,\dd x\right)^{(n-1)/n}
\leq C\int_{\R^n}|\varepsilon(u)(x)|\,\dd x
\]
for compactly supported vector fields, and for an isoperimetric interpretation of this optimal inequality. Here \(\dd x\) denotes Lebesgue measure. The displayed exponent requires an integer dimension \(n\geq2\); the record does not specify a particular dimension or its dimension quantifier.

The conventional meaning of the symmetric gradient is
\[
\varepsilon(u)=\frac{Du+(Du)^{\mathsf T}}2,
\qquad
\varepsilon(u)_{ij}=\frac{\partial_j u_i+\partial_i u_j}{2}
\]
for a differentiable field \(u:\R^n\to\R^n\). This convention, the vector and matrix norms denoted by \(|\cdot|\), and the precise admissible regularity class require confirmation from the source; compact support alone does not specify these choices.

More precisely, once a dimension, norms, and an admissible class \(\mathcal A\) of compactly supported vector fields have been fixed, the requested constant is
\[
C_*(n;\mathcal A)
=\inf\left\{C\geq0:
\|u\|_{L^{n/(n-1)}(\R^n)}
\leq C\int_{\R^n}|\varepsilon(u)|\,\dd x
\text{ for every }u\in\mathcal A\right\},
\]
with value \(+\infty\) if no finite constant works. The task is to determine this constant and an associated isoperimetric interpretation. This is a conditional formulation: the record does not identify the geometric objects, perimeter or energy, or equality correspondence intended by that interpretation.

\textbf{Shape-preserving approximation questions.} The record also asks whether a \(C^2\) approximant to a convex function can be chosen convex, and whether a \(C^2\) approximant to a subharmonic function can be chosen subharmonic.

For clarity, a real-valued function \(f\) on a convex set \(D\subseteq\R^d\) is convex if
\[
f((1-t)x+ty)\leq(1-t)f(x)+tf(y)
\quad(x,y\in D,\ 0\leq t\leq1).
\]
For the real-valued, locally integrable version on an open set \(\Omega\subseteq\R^d\), subharmonicity means upper semicontinuity together with
\[
f(x)\leq\frac{1}{|B(x,r)|}\int_{B(x,r)}f(y)\,\dd y
\quad\text{whenever }\overline{B(x,r)}\subset\Omega,
\]
where \(B(x,r)\) is the Euclidean ball and \(|B(x,r)|\) its Lebesgue measure. These definitions explain the shape conditions; they do not resolve the intended domains or whether extended-valued functions are allowed.

The approximation questions cannot be given complete quantifiers from this record: it supplies neither the approximation relation or error criterion, nor the admissible input class, domains, boundary requirements, or any exceptional-set conditions. In particular, the phrase ``the \(C^2\) approximant'' does not identify an underlying approximation theorem. These data must be recovered before a precise existence assertion for the approximants can be stated.
\subsection{Short English statement}
Determine the sharp endpoint Korn--Sobolev constant for compactly supported vector fields and its isoperimetric meaning. Also ask whether the intended \(C^2\) approximations preserve convexity or subharmonicity. The norm conventions, admissible field class, and especially the approximation criterion are missing, so the supplied record supports only the qualified formulation above. Its partially solved status is supplied dated curation.
\subsection{Sources}
\begin{itemize}
\item Local UnsolvedMath registry record supplied for permanent ID \(30001417\), source problem code \texttt{catalogue entry}, titled ``Sharp Korn--Sobolev Constants and Shape-Preserving Approximation,'' in the Oberwolfach Reports Open Problems collection. The record gives the proposed year as 2010 and names no individual proposer.
\item Supplied report citation: \emph{Optimal Constants in the Theory of Sobolev Spaces and PDEs} (2010), Oberwolfach Reports, DOI 10.4171/owr/2010/08, pp.~325--371. \url{https://doi.org/10.4171/owr/2010/08}. The supplied background describes Oberwolfach Reports as records of mathematical workshops and their open-problem discussions; no precise page or passage for these questions is supplied.
\item Supplied literature reference: Imomkulov, 1993. \url{https://doi.org/10.1070/IM1993v041n01ABEH002184}.
\item Supplied literature reference: Azagra--Drake--Hajlasz, 2023. \url{https://arxiv.org/abs/2311.03481}.
\item Supplied literature reference: Oberwolfach Report 8/2010. \url{https://doi.org/10.4171/owr/2010/08}.
\item Literature assessment dated exactly \texttt{2026-08-21}: the supplied curation labels the record partially solved, reports \(C^2\) approximation results for important convex subclasses, and states that the cited subharmonic theorem does not preserve subharmonicity of the approximant. It describes the full convex-preserving and subharmonic-preserving \(C^2\) questions as unresolved. This is supplied dated curation derived from the OWR material, not an independently verified literature assessment; its time-sensitive status and cited claims require independent verification. The supplied assessment does not specify the status of the sharp Korn--Sobolev constant separately.
\end{itemize}
\end{document}
